% !TEX root = ../main.tex

\begin{tcolorbox}[colback=bluePoli!8, colframe=bluePoli, title=\textbf{Abstract}, arc=2pt]
\small
ABSTRACT
% Vision-Language-Action (VLA) models represent a paradigm shift in robotics and embodied artificial intelligence, coupling visual perception, open-ended language conditioning, and robotic action trajectory generation. However, their physical deployment on edge robotic platforms is hindered by severe hardware constraints in latency, power, and memory bandwidth.

% This report documents the ongoing Multidisciplinary Project on \textit{Workload-Driven Exploration of Edge NPU Architectures for Vision-Language-Action Models}. Rather than prescribing a fixed hardware accelerator upfront, we pursue an empirical, workload-driven methodology:
% \begin{enumerate}[leftmargin=*]
%     \item We characterize the computational anatomy of representative open-source VLA models (principally \texttt{SmolVLA} and $\pi_0$), extracting layer-by-layer GEMM/GEMV dimensions $(M, N, K)$, arithmetic intensity, and tensor dataflows via PyTorch profiling.
%     \item We perform extensive design-space exploration across systolic array dimensions, dataflows, and on-chip SRAM allocations using \textbf{SCALE-Sim v3}.
%     \item We selectively validate critical kernels using cycle-accurate modeling in \textbf{PyTorchSim}.
% \end{enumerate}
% Our primary objective is to investigate whether a homogeneous systolic NPU suffices across all VLA phases, or whether the distinct operational characteristics of the vision backbone and iterative action generation mandate a \textbf{heterogeneous or reconfigurable} edge NPU architecture.
\end{tcolorbox}
\vspace{0.4cm}
